% ---- conclusion.tex ----
% Data sources: /home/sharaths/projects/Prabh\=asa-integration/research/memory/journal.md (closure notes, RQ-C/RQ-D queued)
% /home/sharaths/projects/Prabh\=asa-integration/site/src/data/results.json (model cards, release status)


Prabhāsa is a controlled, pre-registered evaluation of whether Pañinian morphosyntactic structure, operationalized through morpheme-boundary embeddings and role-aware masking, improves data-efficient pretraining of small language models. The experiments isolate architectural effects (RoPE), objective effects (pure MLM), and mechanism effects (kāraka masking and auxiliary supervision) through matched-budget comparisons and multi-seed validation.

\subsection{Summary of Findings}

Three empirical findings emerge from the complete experimental run. The first finding addresses the choice of pretraining objective, the second tests a core Pañinian mechanism, and the third evaluates supervised auxiliary supervision on role prediction.

Regarding the pretraining objective, the objective effect is scale-dependent. Removing the causal language modeling (CLM) head in favor of pure masked-language modeling (MLM) yields no significant difference at 10M parameters, with pure-MLM achieving 63.58 +/- 1.73 versus hybrid 64.09 +/- 0.26 on the BLiMP suite (overlapping 95\% confidence intervals). At 100M parameters, however, pure-MLM substantially outperforms hybrid objectives with a score of 73.06 compared to 67.57, a gain of 5.49 percentage points. This finding motivates the use of pure-MLM training for the Strict-track 100M model while retaining the hybrid objective for the Strict-Small 10M submission. The result suggests that the dilutive effect of the CLM head on MLM gradient flow becomes material at larger scales, consistent with multitask learning literature \cite{Kaplan2020}.

The second finding concerns kāraka-aware masking. A pre-registered controlled comparison between role-stratified masking (Arm K) and uniform random masking (Arm C), with all other factors held constant (budget, architecture, objective, corpus), yields a negligible difference on the targeted agreement and argument-structure paradigms (82.03 versus 81.93, 95\% CI [\(-0.99, +1.20\)]) and on full BLiMP (71.77 versus 70.49). The apparent gains claimed earlier in development were confounded by mask-rate variation and frequency weighting, not role structure itself. This null result provides a well-controlled closure of a pre-registered hypothesis regarding the causal efficacy of morpheme-aware masking.

The third finding addresses auxiliary objectives for role prediction. A pre-registered 5-seed paired $t$-test on the supervised kāraka-role prediction auxiliary objective yields a mean difference of +0.76 percentage points on targeted paradigms (95\% CI [\(-0.74, +2.27\)], $t = 1.41$, $\text{df} = 4$, not statistically significant). The effect attenuated across increasing numbers of seeds (starting at +2.65 for seed 0, reaching +0.76 by the fifth seed), indicating that the initial single-seed result reflected seed-driven variance rather than a robust signal. A shuffled-role control confirmed that the small effect was not attributable to generic multitask regularization. Overall BLiMP performance on the auxiliary objective reached 64.89 compared to baseline 63.88, a +1.01 percentage point difference within the expected noise margin of development.

The primary conclusion is that architecture (RoPE) and objective design (pure MLM at scale) are the primary levers for improving BabyLM performance at the 10M-100M parameter regime. In contrast, the Pañinian mechanisms (morpheme-aware masking and supervised role prediction) do not carry measurable causal weight on standard linguistic diagnostics. This represents a valid scientific outcome derived from pre-registered, multi-seed, well-controlled experiments.

\subsection{Future Work}

Four concrete research directions remain open. The first direction concerns the behavior of Pañinian mechanisms at larger parameter budgets. The scale-dependent objective effect (F1) motivates a full replication of the kāraka auxiliary objective experiment at 100M parameters with extended pretraining (50M or more tokens per seed). At larger scales and with longer training, the model may develop greater capacity to exploit role-structure signals. This experiment is essential for determining whether the non-significance of F3 at 10M reflects a scale artifact or a fundamental property of the auxiliary objective.

A second direction involves evaluating Prabhāsa on benchmarks that probe compositional generalization more directly. Evaluation on agreement and argument-structure diagnostics (BLiMP subsets) may not be the optimal readout for structure-aware pretraining. Compositional generalization benchmarks including SCAN, COGS, and CFQ are better suited to reveal effects of morphosyntactic structure. A targeted replication comparing Pañinian mechanisms to baselines on these benchmarks would test whether the null BLiMP finding generalizes or whether the mechanisms provide benefits on alternative evaluation surfaces.

A third direction explores the application of Prabhāsa to morphologically richer languages. Prabhāsa is evaluated on English-only pretraining, yet the linguistic design of Pañinian kāraka roles suggests that effects may be larger in morphologically rich languages such as Sanskrit, Finnish, or Turkish, or in zero-shot cross-lingual transfer settings. A multilingual variant using Sanskrit pre-pretraining with structured role annotations from the Sanskrit Heritage Lexicon \cite{Huet2003SanskritHeritage}, followed by English fine-tuning, would test whether structure-aware mechanisms show stronger signatures when applied to syntactically complex source languages.

The fourth direction addresses post-training enhancements and reasoning quality. The original Prabhāsa program includes a Navya-Nyāya logical reasoning scaffold designed to improve inferential validity via a 6-phase Pañcāvayava fine-tuning pipeline. This component is out of scope for the BabyLM track but represents a distinct evaluation surface where Pañinian structure may provide benefits beyond standard supervised benchmarks. A targeted experiment on fallacy detection and inference-quality metrics (tested on a small annotated corpus of logical-validity judgments or on the subset of BLiMP that probes inferential consistency) would measure whether the best-performing H1 arm (pure-MLM at 100M) synergizes with the reasoning scaffold.

\subsection{Limitations}

The primary limitation is the small parameter range explored (10M to 100M parameters). While the BabyLM budget is by design constrained, effects that emerge only at larger scales beyond 350M parameters would not be detected in this study. The scale-dependent objective effect (F1) demonstrates that large parameter differences can substantially change mechanism efficacy; the current study is therefore underpowered to rule out structure-aware mechanisms at larger parameter scales.

A second limitation is the restricted evaluation scope. BLiMP focuses on agreement and argument structure; compositional generalization, semantic role labeling, and cross-lingual transfer are not measured. The null findings on BLiMP do not rule out effects on these alternative evaluation surfaces.

A third limitation concerns the single-seed nature of the Strict-track submission model (100M pure-MLM, BLiMP 73.06). While the objective-effect comparison (F1) uses 3 seeds and the kāraka effects (F2, F3) use either 1 seed per arm or 5 seeds, the final submitted model represents a single instantiation of the best arm. Replication with multiple seeds would strengthen the Strict-track baseline claims and provide tighter confidence bounds on the reported performance.

\subsection{Reproducibility and Release}

All code, data, and results are released under open-source license. The GitHub repository at \texttt{github.com/SharathSPhD/prabhasa-babylm} contains 269 commits from a single contributor and captures the full development history. Trained models are published on HuggingFace: \texttt{qbz506/prabhasa-b\_s} for the 100M Strict-track model and \texttt{qbz506/prabhasa-b\_ss-0.1} for the Strict-Small submission model. Each model includes a detailed card disclosing methods, evaluation protocol, and explicit limitations. The dataset \texttt{qbz506/prabhasa-babylm-grammar} provides morpheme inventories and kāraka role labels for corpus preprocessing, enabling external reproduction of the pre-processing pipeline. Intermediate checkpoints, per-seed BLiMP result tables, Tarka adversarial review memos, and the complete experiment ledger are archived in \texttt{research/memory/} and \texttt{docs/decisions/} within the repository.

The experiment ledger records the attempt number, configuration changes, results, and interpretation for each run, enabling full transparency and external audit of the development process and the closure decisions. This documentation supports independent verification of the findings and encourages extensions of the work to new evaluation surfaces and parameter regimes.
